\documentclass[10pt,a4paper]{amsart}
\usepackage[margin=19mm]{geometry}
\usepackage{amsmath,amssymb,mathtools}
\usepackage[scr=rsfs,cal=euler]{mathalfa}
\usepackage{xfrac}
\usepackage[unicode,hidelinks,bookmarks=false]{hyperref}
\newcommand{\ie}{i.e.\ }
\newcommand{\cf}{cf.\ }
\newcommand{\resp}{respectively\ }
\newcommand{\Subsection}{\subsection}
\providecommand{\nobreakdash}{\nobreak\hbox{-}}
\setlength{\emergencystretch}{2em}
\multlinegap0pt
\usepackage{bm}
\usepackage{bbm}
\usepackage{dsfont}
\usepackage{xspace}
\usepackage{cancel}
\usepackage{mathtools}
\usepackage{amsthm}
\makeatletter
\@ifundefined{lemma}{\newtheorem{lemma}{Lemma}}{}
\makeatother
\title[A scheme-theoretic interpretation of generalized splines]{A scheme-theoretic interpretation of generalized splines}
\author{Kyle Stoltz}
\date{Source: arXiv:2509.05840v3 (CC BY 4.0)}
\begin{document}
\maketitle

\noindent\emph{Source and licence.} A scheme-theoretic interpretation of generalized splines. Version arXiv:2509.05840v3 (CC BY 4.0), distributed under the Creative Commons Attribution 4.0 International licence, as recorded in the arXiv licence metadata for this exact version and shown on the abstract page https://arxiv.org/abs/2509.05840v3. Full rights notice: licenses/arxiv-2509.05840-CC-BY-4.0.txt.\par\smallskip

\noindent\textit{Editorial scope.} Counted unit: the Lemma whose source label is \texttt{lma:graph-pushforward} in the article (source0.tex lines 503--506, source label \texttt{lma:graph-pushforward}), delivered together with the article's own complete original proof of it (source0.tex lines 507--532, one \texttt{proof} environment of 26 lines). No mathematics is written, repaired or re-derived: statement and proof are the article's own text line for line. Required context retained uncounted: none. Every definition, display and label of those lines is retained unchanged. Omitted from this package and not redistributed: 1--502, 533--1873. Nothing inside a retained excerpt is omitted or altered: every retained body is delivered exactly as the article's lines are written, so no omission receipt of any kind accompanies this package. Citations: the retained excerpts carry no citation command at all, so no bibliography entry of the article is transcribed and no reference is redistributed. Reference adaptation: none was needed and none was made; every reference inside the delivered unit resolves inside it, so no unresolved reference or citation is produced. Printed slips kept literal: no printed word, symbol, hypothesis or number of the counted statement or of its proof was altered. Layout and class: the article's own class is \texttt{article} with packages including graphicx, tikz, tikz-cd, amsmath, amsfonts, amsthm, amssymb, tabularx, booktabs, url; the wrapper is the lane's pinned \texttt{amsart} editorial wrapper at 10pt on A4, which restates the theorem-like environments the retained text needs and adds the article's own macro definitions that the retained text needs (none), with extra wrapper packages (bm, bbm, dsfont, xspace, cancel, mathtools). The delivered numbering is the unit's own. Assets: no retained span contains a figure environment, a graphics-inclusion command, a PDF-page inclusion, a table or a TikZ picture; no image file is copied into this package, the package declares no asset dependency and no image file is redistributed with it. Licence: version arXiv:2509.05840v3 of this article is distributed under the Creative Commons Attribution 4.0 International licence, as recorded in the arXiv licence metadata for this exact version and shown on the abstract page https://arxiv.org/abs/2509.05840v3. A full rights notice is delivered at licenses/arxiv-2509.05840-CC-BY-4.0.txt. Characters: every retained excerpt is pure ASCII, so the delivered engine, XeLaTeX with its default Latin Modern text font, reports no missing character.
\begin{lemma}
\label{lma:graph-pushforward}
Let $R$ be a commutative ring, $G$ an edge-labeled graph, and $\mathfrak G$ the associated spline diagram of rings. Suppose $F_S(-) = - \otimes_R S$ with $S$ flat. Then there exists an edge-labeled graph denoted $F_S(G)$ over $F_S(R)$ with associated spline diagram of rings $F_S(\mathfrak G)$.
\end{lemma}

\begin{proof}
Let $R$, $G$, $\mathfrak G$, and $F_S$ be given as specified. Since $F_S$ is a functor we know $F_S(\mathfrak G)$ acts on edge diagrams as follows:

\[
R_u \rightarrow R_u/I_{uv} \leftarrow R_v
\]
\[
F_S(R_u) \rightarrow F_S(R/I_{uv}) \leftarrow F_S(R_v).
\]
Now in particular, since $F_S$ is given by tensor product with $S$ flat we have the following exact sequences of $R$-modules

\[
0 \rightarrow I_{uv} \rightarrow R_u \rightarrow R/I_{uv} \rightarrow 0
\]
\[
0 \rightarrow F_S(I_{uv}) \rightarrow F_S(R_u) \rightarrow F_S(R/I_{uv}) \rightarrow 0.
\]

Hence each edge diagram under $F_S$ is of the form

\[
F_S(R_u) \rightarrow F_S(R)/F_S(I_{uv}) \leftarrow F_S(R_v) 
\]

Now let $G$ be the original edge-labeled graph, and take $R$ to $F_S(R)$ and $F_S(G)$ to the edge-labeled graph obtained by replacing each edge ideal $I_{uv}$ with $F(I_{uv})$. Then by the previous observation, we can see $F_S(G)$ has $F_S(\mathfrak G)$ as its associated ring of limit splines. 
\end{proof}

\end{document}
